\section{Orientability from two sparse histograms}
\label{sec:inversion}

\subsection{A general two-sheeted-cover identity}

Let an orbit relation carry additive weights in $\Z^d$ and let a
two-sheeted parity cover of the relation be specified. Lift each point's
weight unchanged to both of its lifts. A component with trivial parity
monodromy lifts to two components with its original weight $v$. A component
with nontrivial parity monodromy lifts to one component of weight $2v$.

Write $H(v)$ for the base histogram, $K(v)$ for the cover histogram, and
$O(v),A(v)$ for the trivial- and nontrivial-monodromy parts of $H$.
The letter $A$ will denote nonorientable surface components in the normal
application. Then
\begin{align}
 H(v)&=O(v)+A(v),\\
 K(v)&=2O(v)+A(v/2).
 \label{eq:cover}
\end{align}
Here $A(v/2)=0$ if $v$ has an odd coordinate, or if the integral half does
not belong to the base support. All histograms have finite support.

\begin{theorem}[Sparse doubling inversion]
\label{thm:inversion}
The pair $(H,K)$ uniquely determines finite-support functions $O,A$ that
satisfy \eqref{eq:cover}, whenever such nonnegative integer functions exist.
For $v\ne0$, process base keys in increasing $\ell^1$ norm and set
\begin{equation}
 A(v)=\frac{2H(v)+A(v/2)-K(v)}{2},\qquad O(v)=H(v)-A(v).
 \label{eq:inversion}
\end{equation}
At zero use
\begin{equation}
 A(0)=2H(0)-K(0),\qquad O(0)=K(0)-H(0).
 \label{eq:zero}
\end{equation}
For fixed $d$ and $s=|\supp H|+|\supp K|$, an implementation using
balanced ordered maps requires
$O(s\log(s+1))$ comparisons and $O(s)$ integer arithmetic operations,
including complete forward verification of the resulting histograms.
\end{theorem}
\begin{proof}
Eliminate $O(v)$ from \eqref{eq:cover}. A nonzero integral half has strictly
smaller norm, so its $A$ value is already known if it occurs in the base
support; otherwise it is zero because $0\le A\le H$. This gives
\eqref{eq:inversion}. At zero the half is zero itself, and solving the two
linear equations gives \eqref{eq:zero}. These arguments prove uniqueness.

For existence, reject a nonintegral numerator or a value outside
$0\le A(v)\le H(v)$. Construct $O=H-A$, then reconstruct the entire cover
histogram by adding $2O(v)$ at $v$ and $A(v)$ at $2v$. The reconstructed
base and cover histograms must agree with both inputs on their full
supports. This final comparison rejects extraneous cover-only keys and
missing doubled contributions. Thus a successful result satisfies exactly
the required equations.

Only actual base keys are processed. No intermediate missing values along
a doubling chain need to be generated. Sorting costs the stated number of
comparisons, and each key contributes a constant number of dictionary
lookups and arithmetic operations. Balanced ordered maps give a
deterministic comparison implementation with the same conservative
logarithmic overhead. These bounds count vector operations at fixed $d$;
if $d$ varies, reading and comparing a signature also charges its dimension.
\end{proof}

For fixed $d$, with each coordinate and multiplicity, including temporary
sums, of at most $\ell$ bits, the ordered-map bit bound is
$O(s\ell\log(s+1))$, apart from container allocation costs. If input
fields have at most $L$ bits, $\ell=L+O(\log(s+1))$ suffices even for
inconsistent histograms rejected by the final check.
The algorithm does not depend on the numerical length of a sparse gap
between two present signatures.

The delivered Python code uses dictionaries. Its usual lookup behaviour
is efficient, but adversarial integer-key hash collisions preclude claiming
the ordered-map bound as its deterministic worst-case runtime. A
conservative dictionary bound is $O(s^2\ell+s\ell\log(s+1))$ here.
The overall supplied-vector census remains polynomial; the sharper sparse
bound describes the stated comparison-map implementation of the same
recurrence. The row-merging estimate below uses the same ordered-map
convention.

\begin{example}[A collision between an annulus and a doubled M\"obius band]
Suppose a surface has three M\"obius bands and five annuli. Its base keys are
$H(0,1)=3$ and $H(0,2)=5$. The cover has thirteen annuli, so
$K(0,2)=13$ and $K(0,1)=0$. Equation \eqref{eq:inversion} first gives
$A(0,1)=3$, and then
$A(0,2)=(10+3-13)/2=0$. Comparing $K(v)$ to $2H(v)$ independently at
each key would miss the contribution arriving from the half-signature.
\end{example}

\begin{example}[The zero signature cannot be treated recursively]
A torus and a Klein bottle both have $(\chi,b)=(0,0)$. If there are $a$
tori and $c$ Klein bottles, then $H(0)=a+c$ and $K(0)=2a+c$.
Equation \eqref{eq:zero} recovers them. Substituting $A(v/2)=0$ at zero
would generally give a false answer. The geometric test corpus includes
a closed normal Klein bottle and its multiples in an orientable manifold
with one torus boundary.
\end{example}

\subsection{Why torsion coordinates need a different formula}

The displayed inversion relies on the uniqueness of an integral half in
$\Z^d$ and on strict norm decrease away from zero. If weight coordinates
are reduced modulo two beforehand, doubling is no longer injective; the
incoming term must sum over all preimages, and the same recurrence is
invalid. This is why the delivered construction keeps exactly the integral
pair $(\chi,b)$. A future extension can retain further \emph{integer}
weights, invert in $\Z^d$, and only then reduce a display field modulo two.
An extension with torsion weights requires a separately proved inversion.

\section{Normal orientation doubles and the topology census}
\label{sec:census}

\begin{lemma}[Normal double and boundary lifts]
\label{lem:normal-double}
Let $F(x)$ be a compact properly embedded normal surface in an orientable
three-manifold. The surface $F(2x)$ is the normal orientation double of
$F(x)$: an orientable component contributes two copies, and a
nonorientable component contributes its connected orientation double.
Every boundary circle has two distinct lifts. In the maintained edge-point
ordering, the projection from doubled to original points is
$q\mapsto\lfloor q/2\rfloor$.
\end{lemma}
\begin{proof}
Take a small normal interval neighbourhood of the surface. Its horizontal
boundary has two local sheets over every normal disc, so its normal
coordinates are $2x$. The sheet transition across a face either preserves
or reverses order. Trivial normal monodromy gives two global sheets;
nontrivial monodromy exchanges them and gives a connected cover. In an
orientable ambient three-manifold, orientability of a surface is equivalent
to orientability of its normal line bundle, so this is precisely the
orientation double. This is the classical normal-coordinate doubling
principle used by AHT \cite{aht}.

The restriction of the orientation character to a boundary circle is
trivial. Indeed its tangent line is orientable and the inward collar line
inside the surface is trivial. Hence that circle has two distinct lifts,
even on a nonorientable surface.

For the explicit point indexing, each original point $p$ is replaced by
the adjacent points $2p,2p+1$. A scaled translation preserves the second
coordinate; a scaled reflection reverses it. Global edge-block offsets
and lengths also double. Therefore the projection is the asserted floor
map, and a base interval $[a,b)$ has full inverse image $[2a,2b)$.
\end{proof}

\begin{lemma}[One boundary transversal suffices]
\label{lem:one-boundary}
Lift the two weights \eqref{eq:two-weights} to $F(2x)$ by replacing every
weight interval $[a,b)$ by $[2a,2b)$ and keeping its vector value. The
weighted sum on each doubled component is its actual pair $(\chi,b)$.
No boundary-orbit discovery on $\partial F(2x)$ is needed.
\end{lemma}
\begin{proof}
Every owned cell of the original Euler weighting lifts to two owned cells,
with one charge on the corresponding owned point of each lift. Thus the
first coordinate is Euler characteristic on each doubled component.
For a boundary circle, its chosen base representative has one lift on
each of the two lifted circles from \cref{lem:normal-double}. Accordingly
the second coordinate has exactly one marker per lifted boundary circle.
The conclusion is componentwise, including when both lifts lie in one
connected orientation double.
\end{proof}

\begin{theorem}[Two-weight connected-topology census]
\label{thm:census}
For an admissible normal vector $x$ in the supported ambient manifold,
the entire type spectrum $\Spec(x)$ can be computed by the following
operations:
\begin{enumerate}
\item one boundary orbit discovery and its least-transversal compilation;
\item one two-weight orbit histogram on $F(x)$;
\item one two-weight orbit histogram on $F(2x)$, using lifted original weights;
\item sparse inversion of the two histograms and surface classification.
\end{enumerate}
The algorithm is polynomial in the binary encoded input size and never
lists individual normal discs, points, or components. Its output rows have
exact binary multiplicities and determine every component homeomorphism
type, grouped by that type.
\end{theorem}
\begin{proof}
\Cref{thm:transversal,thm:nested} and the Euler ownership construction
identify the first histogram with $(\chi,b)$ on each component.
\Cref{lem:normal-double,lem:one-boundary} identify the second histogram
with the normal orientation-cover histogram and establish
\eqref{eq:cover} with $O$ orientable and $A$ nonorientable components.
\Cref{thm:inversion} recovers their separate multiplicities.
The classification of compact connected surfaces then determines the
genus or crosscap number from $(\chi,b)$ and orientability. The
implementation rejects impossible or nonintegral classification data.

The normal-arc relations have $O(t)$ pairings and
$O(B+\log(t+1))$-bit endpoints. The boundary selector has at most
$G_\partial$ intervals; after its inclusion in the full universe there
are $O(t+G_\partial)$ weight intervals. The inherited AHT bounds and
\cref{thm:transversal} make all these quantities polynomial in $t,B$.
Weighted AHT in fixed dimension and the sparse final inversion preserve
that bound.
\end{proof}

\begin{center}
\begin{tabular}{lrrcl}
\toprule
Connected type & $\chi$ & $b$ & Orientable? & Recovered parameter\\
\midrule
Sphere & 2 & 0 & yes & $g=0$\\
Disc & 1 & 1 & yes & $g=0$\\
Annulus & 0 & 2 & yes & $g=0$\\
M\"obius band & 0 & 1 & no & $k=1$\\
Torus & 0 & 0 & yes & $g=1$\\
Klein bottle & 0 & 0 & no & $k=2$\\
Projective plane & 1 & 0 & no & $k=1$\\
Genus $g$, $b$ boundary circles & $2-2g-b$ & $b$ & yes & $g$\\
$k$ crosscaps, $b$ boundary circles & $2-k-b$ & $b$ & no & $k$\\
\bottomrule
\end{tabular}
\end{center}

The projective-plane row is retained because the ambient contract does not
assume irreducibility. Deleting it as an alleged impossibility would make
the implementation incorrect on a supported fixture.
